%% guide-kga.tex　ガイドブック2冊（子ども用 guide-kids.tex・大人用 guide-adult.tex）で共通の部品（組 G-A、頭文字 kga）
%% 第16回編集委員会（10月30日 10:14 の依頼者の指示で2冊に分けた）。親ファイルの前付けと奥付・子ども用の本文で使う。
\makeatletter
%% 色：kgaPA＝子ども用の札の色（れんが色。白い地で約5.9:1）、kgaPB＝大人用の札の色（巻の色 bkVol0）
%% 札は色だけでなく、字（子ども用／大人用）と形（丸／四角）でも区別する
\definecolor{kgaPA}{cmyk}{0.05,0.65,0.95,0.28}
\colorlet{kgaPB}{bkVol0}
\newcommand\kgaNoRuby{\def\R##1##2{##1}}% ルビを外す
%% 本ごとに親ファイルで決めるもの（書名・副題・キャッチ・帯の小さな字・偶数ページの柱）
\providecommand\kgaBookTitle{ガイドブック}
\providecommand\kgaBookSub{}
\providecommand\kgaCatch{}
\providecommand\kgaBandLabel{}
%% ラベルのページ番号（まだ組んでいない所は「—」）  \kgaPage{ラベル}
\newcommand\kgaPage[1]{\@ifundefined{r@#1}{—}{\pageref{#1}}}

%% --------------------------------------------------------------------
%% 9冊の本のたなの絵  \kgaShelf{x}{y}{倍率}{棚板の色}：本の高さはページ数に比例（400〜640ページ）
%% --------------------------------------------------------------------
\newcommand\kgaShelf[4]{%
  \begin{scope}[shift={(#1,#2)}, scale=#3]
    \foreach \v/\p in {1/400,2/456,3/564,4/560,5/492,6/600,7/492,8/632,9/640} {
      \pgfmathsetmacro\kga@bx{(\v-1)*16.2}\pgfmathsetmacro\kga@bh{\p*0.13}
      \fill[bkVol\v, rounded corners=0.8mm] (\kga@bx,0) rectangle ++(14.4,\kga@bh);
      \draw[bkBandOn\v, line width=0.5pt, opacity=0.55] (\kga@bx+1.6,\kga@bh-6) -- ++(11.2,0) (\kga@bx+1.6,6) -- ++(11.2,0);
      \node[text=bkBandOn\v, font=\sffamily\bfseries\fontsize{15}{15}\selectfont] at (\kga@bx+7.2,\kga@bh-14) {\v};
    }
    \fill[#4, rounded corners=0.6mm] (-3,-3.2) rectangle (148.2,0);
  \end{scope}}

%% --------------------------------------------------------------------
%% 表紙の札  \kgaBadge{kids|adult}{中心x}{中心y}{大きさ（1＝表紙）}
%%   子ども用：白い丸に、れんが色の太い輪と「子ども用」。大人用：白い四角に、紺灰色の二重の枠と「大人用」
%% --------------------------------------------------------------------
\newcommand\kgaBadge[4]{%
  \begin{scope}[shift={(#2,#3)}, scale=#4, transform shape]
  \def\kga@b{#1}\def\kga@kids{kids}%
  \ifx\kga@b\kga@kids
    \fill[white] (0,0) circle (14);
    \draw[kgaPA, line width=2.4pt] (0,0) circle (12.6);
    \node[text=kgaPA, font=\sffamily\bfseries\fontsize{15}{17}\selectfont, inner sep=0pt] at (0,0) {\R{子}{こ}ども\R{用}{よう}};
  \else
    \fill[white] (-17,-11) rectangle (17,11);
    \draw[kgaPB, line width=1.6pt] (-15.4,-9.4) rectangle (15.4,9.4);
    \draw[kgaPB, line width=0.6pt] (-13.6,-7.6) rectangle (13.6,7.6);
    \node[text=kgaPB, font=\sffamily\bfseries\fontsize{14}{16}\selectfont, inner sep=0pt] at (0,0) {大人用};
  \fi
  \end{scope}}

%% --------------------------------------------------------------------
%% 表1（表紙）：\bkFrontArt と同じ配置。巻の番号の札のかわりに「子ども用」「大人用」の札
%%   \kgaFrontArt{kids|adult}
%% --------------------------------------------------------------------
\newcommand\kgaFrontArt[1]{%
  \fill[bkVol] (-3,64) rectangle (185,260);
  \fill[white] (-3,-3) rectangle (185,64);
  \node[anchor=base west, inner sep=0pt, text=bkOn, font=\sffamily\bfseries\fontsize{12.5}{15}\selectfont] at (16,242) {\bkSeries};
  \draw[bkOn, line width=0.6pt] (16,237.5) -- (112,237.5);
  \kgaBadge{#1}{147}{236}{1}
  \node[anchor=base west, inner sep=0pt, text=bkOn, font=\bk@titlefont] at (15,206) {\kgaBookTitle};
  \node[anchor=base west, inner sep=0pt, text=bkOn, font=\sffamily\bfseries\fontsize{17}{20}\selectfont] at (16,192) {\kgaBookSub};
  \node[anchor=base west, inner sep=0pt, text=bkOn, font=\sffamily\fontsize{11.5}{14}\selectfont] at (16,181.5) {\kgaCatch};
  \foreach \r in {1,...,7} {\node[inner sep=0pt] at ({20+(\r-1)*11.5},164) {\bkrouteicon[7.5mm]{\r}};}
  \kgaShelf{18}{74}{1}{white}
  \bk@lineup{16}{36}
  \node[anchor=base, text=bkGray, font=\sffamily\fontsize{8}{10}\selectfont] at (91,16) {\R{全}{ぜん}9\R{巻}{かん}　\bkSeries};
  \node[anchor=base, text=bkGray, font=\sffamily\fontsize{7.5}{9}\selectfont] at (91,8.5) {\bkEditor\hspace{1pt}\R{編}{へん}};
}

%% --------------------------------------------------------------------
%% 表4（裏表紙）  \kgaBackArt{見出し}{中身（TikZ。y=200 から下へ）}{ふきだしの文}
%% --------------------------------------------------------------------
\newcommand\kgaBackArt[3]{%
  \fill[bkVol!9] (-3,-3) rectangle (185,260);
  \fill[bkVol] (-3,226) rectangle (185,260);
  \node[anchor=base west, inner sep=0pt, text=bkOn, font=\sffamily\fontsize{9.5}{12}\selectfont] at (16,246) {\bkSeries};
  \node[anchor=base west, inner sep=0pt, text=bkOn, font=\sffamily\bfseries\fontsize{19}{24}\selectfont] at (16,233)
    {\kgaBookTitle\hspace{0.5em}{\fontsize{12}{14}\selectfont\kgaBookSub}};
  \node[anchor=base west, inner sep=0pt, text=bkDeep, font=\sffamily\bfseries\fontsize{13}{16}\selectfont] at (16,210) {#1};
  \draw[bkVol, line width=0.8pt] (16,206) -- (166,206);
  #2
  \node[anchor=base west, inner sep=0pt, text=bkDeep, font=\sffamily\bfseries\fontsize{13}{16}\selectfont] at (16,140) {\R{全}{ぜん}9\R{巻}{かん}のならび};
  \draw[bkVol, line width=0.8pt] (16,136) -- (98,136);
  \foreach \i in {1,...,9} {
    \pgfmathsetmacro\bk@ly{127-(\i-1)*8.6}
    \fill[bkVol\i, rounded corners=0.8mm] (16,\bk@ly-2.8) rectangle ++(5.6,5.6);
    \node[text=bkBandOn\i, font=\sffamily\bfseries\fontsize{8}{8}\selectfont] at (18.8,\bk@ly) {\i};
    \node[anchor=west, inner sep=0pt, text=bkInk, font=\sffamily\fontsize{9.5}{11}\selectfont] at (24.5,\bk@ly-0.3)
      {\ltjsetruby{stretch=001}\bkTitle{\i}};
  }
  \node[rectangle callout, callout absolute pointer={(136.5,91)}, anchor=south, draw=bkVol, line width=0.8pt, fill=white,
    rounded corners=3mm, inner xsep=3.5mm, inner ysep=2.5mm, text width=50mm, align=flush center, text=bkDeep,
    font=\sffamily\bfseries\fontsize{10.5}{16}\selectfont] at (136.5,97) {#3};
  \fill[bkVol!22, rounded corners=2mm] (107,50) rectangle (166,54);
  \bkChar{sora}{(117,53.5)}{31mm}\bkChar{mina}{(136.5,53.5)}{31mm}\bkChar{delta}{(156,53.5)}{19mm}
  \draw[bkGray, line width=0.6pt, dash pattern=on 2pt off 1.5pt, fill=white] (114,12) rectangle (166,38);
  \node[text=bkGray, align=flush center, font=\sffamily\fontsize{7}{9}\selectfont] at (140,25) {ISBN・バーコード\\（入稿時に配置）};
  \node[anchor=base west, inner sep=0pt, text=bkGray, font=\sffamily\fontsize{7.5}{10}\selectfont] at (16,25) {\bkSeries　\kgaBookTitle};
  \node[anchor=base west, inner sep=0pt, text=bkGray, font=\sffamily\fontsize{7.5}{10}\selectfont] at (16,16) {\bkEditor\hspace{1pt}\R{編}{へん}};
}
% 表4の中身の1行（札つき）  \kgaBackItem{y}{札の文字}{文}
\newcommand\kgaBackItem[3]{%
  \node[anchor=west, fill=bkVol, text=bkOn, rounded corners=1.2mm, inner xsep=1.6mm, inner ysep=1mm, minimum width=7mm,
    font=\sffamily\bfseries\fontsize{8.5}{10}\selectfont] (kgabi) at (16,#1) {#2};
  \node[anchor=west, inner sep=0pt, text=bkInk, font=\sffamily\fontsize{10}{13}\selectfont] at ([xshift=2.5mm]kgabi.east) {#3};}

%% --------------------------------------------------------------------
%% とびら（p.1）  \kgaTitlePage{kids|adult}
%% --------------------------------------------------------------------
\newcommand\kgaTitleArt[1]{%
  \fill[bkVol] (-3,246) rectangle (185,260);
  \node[anchor=base, text=bkDeep, font=\sffamily\bfseries\fontsize{12}{15}\selectfont] at (91,206) {\bkSeries};
  \kgaBadge{#1}{91}{188}{0.85}
  \node[anchor=base, text=bkDeep, font=\sffamily\bfseries\fontsize{32}{40}\selectfont] at (91,158) {\kgaBookTitle};
  \node[anchor=base, text=bkDeep, font=\sffamily\bfseries\fontsize{15}{18}\selectfont] at (91,145) {\kgaBookSub};
  \node[anchor=base, text=bkGray, font=\fontsize{11}{15}\selectfont] at (91,130) {\kgaCatch};
  \kgaShelf{52.5}{78}{0.5}{bkVol}
  \node[anchor=base, text=bkGray, font=\sffamily\fontsize{9}{12}\selectfont] at (91,34) {\bkEditor\hspace{1pt}\R{編}{へん}};
  \fill[bkVol] (-3,-3) rectangle (185,4);
}
\newcommand\kgaTitlePage[1]{%
  \clearpage\setcounter{page}{1}%
  \begingroup\global\bk@decofalse\bkfolio{false}%
  \AddToHookNext{shipout/background}{\put(0,-\paperheight){%
    \begin{tikzpicture}[overlay, x=1mm, y=1mm]\kgaTitleArt{#1}\end{tikzpicture}}}%
  \null\clearpage\endgroup\global\bk@decotrue\bkfolio{true}}

%% --------------------------------------------------------------------
%% もくじ
%% --------------------------------------------------------------------
% 行  \kgaTocRow{色}{番号}{題名}{ひとこと（空なら1行）}{ラベル}
\newcommand\kgaTocRow[5]{%
  \par\noindent
  \parbox[t]{12mm}{\leavevmode\tikz[baseline=(p.base)]\node[fill=#1, text=white, rounded corners=1.3mm,
      minimum width=9mm, inner xsep=1mm, inner ysep=1.1mm, font=\sffamily\bfseries\fontsize{9.5}{10}\selectfont] (p) {#2};}%
  \parbox[t]{\dimexpr\linewidth-12mm-16mm\relax}{\setlength\parskip{0pt}%
    \leavevmode{\sffamily\bfseries\fontsize{11}{17}\selectfont\color{bkInk}#3}\if\relax\detokenize{#4}\relax\else\par\nobreak\vspace{-0.6mm}
    {\rmfamily\fontsize{9.5}{15}\selectfont\color{bkInk}#4}\fi}%
  \parbox[t]{16mm}{\raggedleft\leavevmode{\sffamily\bfseries\fontsize{11}{17}\selectfont\color{bkInk}\bkpageno{\kgaPage{#5}}}}%
  \par\vspace{1.2mm}{\color{bkLine}\hrule height 0.35pt}\vspace{1.2mm}}
% 見出し  \kgaTocHead
\newcommand\kgaTocHead{%
  \noindent\begin{tikzpicture}
    \node[anchor=base west, inner sep=0pt, text=bkDeep, font=\sffamily\bfseries\fontsize{30}{36}\selectfont] (h) at (0,0) {もくじ};
    \node[anchor=base west, inner sep=0pt, text=bkDeep, font=\sffamily\bfseries\fontsize{12}{15}\selectfont] at ([xshift=6mm]h.base east) {\kgaBookTitle\hspace{0.6em}\kgaBookSub};
    \draw[bkVol, line width=1.6pt, overlay] (0,-3.5mm) -- (\linewidth,-3.5mm);
  \end{tikzpicture}\par\vspace{6mm}}
% \kgaTocPages{1ページめ}{2ページめ（空ならもくじは1ページ）}
\newcommand\kgaTocPages[2]{%
  \clearpage\setcounter{page}{2}%
  \global\bk@decofalse\bkfolio{true}\bksetrunning{}{}{}%
  \kgaTocHead #1\par\clearpage
  \if\relax\detokenize{#2}\relax\else #2\par\clearpage\fi
  \global\bk@decotrue}

%% --------------------------------------------------------------------
%% 色の面の上に問いを置いたページ（子ども用の「はじめに」、p.3）。飾り・ノンブルなし
%%   \kgaArtPage{大きな字}{ラベル}{題名}{問い y/文,…}{絵}{人物}{ふきだしの文}
%% --------------------------------------------------------------------
\newcommand\kgaArtPage[7]{%
  \clearpage
  \begingroup\global\bk@decofalse\bkfolio{false}%
  \AddToHookNext{shipout/background}{\put(0,-\paperheight){%
  \begin{tikzpicture}[overlay, x=1mm, y=1mm]
    \fill[bkVol] (-3,104) rectangle (185,260);
    \begin{scope}\clip (-3,104) rectangle (185,260);
      \node[text=white, opacity=0.12, font=\sffamily\bfseries\fontsize{200}{200}\selectfont] at (140,190) {#1};
    \end{scope}
    \node[anchor=base west, text=white, font=\sffamily\bfseries\fontsize{16}{20}\selectfont] at (22,230) {#2};
    \node[anchor=north west, text=white, font=\sffamily\bfseries\fontsize{30}{40}\selectfont, text width=140mm, align=left, inner sep=0pt] at (21,222) {#3};
    \foreach \y/\t in {#4} {
      \node[anchor=west, inner sep=0pt] at (22,\y) {\bkicon[1.6em]{hatena}{bkYellow}};
      \node[anchor=west, text=white, font=\sffamily\bfseries\fontsize{12.5}{16}\selectfont] at (31,\y) {\t};
    }
    #5
    \bkChar{#6}{(28,62)}{30mm}
    \node[rectangle callout, callout absolute pointer={(43,80)}, anchor=west, draw=bkVol, line width=0.8pt, fill=white,
      rounded corners=3mm, inner xsep=4mm, inner ysep=3mm,
      text width=94mm, align=left, font=\fontsize{10.5}{18}\selectfont, text=bkInk] at (49,80) {#7};
  \end{tikzpicture}}}%
  \null\clearpage\endgroup\global\bk@decotrue\bkfolio{true}}

%% --------------------------------------------------------------------
%% 子ども用の本文の節の見出し（偶数ページから）  \kgasec{番号}{題名}{ひとこと}{しおり用の文字}
%%   ラベル kga-sec番号 を置く（もくじが拾う）。柱：偶数ページ＝\kgaBandLabel、奇数ページ＝番号と題名
%% --------------------------------------------------------------------
\newcommand\kgasec[4]{%
  \bkevenpage
  \global\bk@decotrue
  \bksetrunning{}{\kgaBandLabel}{#1\hspace{0.6\zw}#2}%
  \pdfbookmark[0]{#1　#4}{kgasec#1}\typeout{BKPAGE kga #1 \thepage}%
  \label{kga-sec#1}%
  \noindent\begin{tikzpicture}
    \useasboundingbox (0,0) rectangle (\linewidth,3.1);
    \begin{scope}\clip[rounded corners=2.5mm] (0,0) rectangle (\linewidth,3.1);
      \fill[bkTheme] (0,0) rectangle (\linewidth,3.1);
      \fill[white, opacity=0.13] (\linewidth-2.2,3.1) circle (2.0);
      \fill[white, opacity=0.10] (\linewidth-0.6,0) circle (1.3);
    \end{scope}
    \node[anchor=west, text=bkOn, font=\sffamily\small] at (0.5,2.6) {\kgaBandLabel};
    \node[anchor=west, circle, fill=white, text=\bkTT, minimum size=1.5cm, font=\sffamily\bfseries\Huge, inner sep=0pt] (n) at (0.45,1.25) {#1};
    \node[anchor=west, text=bkOn, font=\sffamily\bfseries\LARGE, text width=\linewidth-3.6cm, align=left, inner sep=0pt] at ([xshift=0.35cm]n.east) {#2};
  \end{tikzpicture}\par\vspace{2.5mm}
  \noindent{\sffamily\fontsize{10.5}{19}\selectfont\color{bkInk}#3}\par
  \vspace{1.5mm}{\color{bkLine}\hrule height 0.5pt}\vspace{3mm}}

%% --------------------------------------------------------------------
%% 奥付  \kgaColophon{書名の行}{副題の行}{版の行}{いちばん下の注}
%% --------------------------------------------------------------------
\newcommand\kgaColophon[4]{%
  \bkpadfour
  \global\bk@decofalse\bkfolio{false}%
  \null\vfill
  \noindent\begin{minipage}{\linewidth}\small
  {\sffamily\bfseries\Large\color{bkVol0}#1}\par\vspace{1mm}
  {\sffamily\bfseries\large #2}\par\vspace{3mm}
  {\color{bkVol0}\hrule height 1pt}\vspace{3mm}
  \begin{tabular}{@{}l@{\hspace{6mm}}l@{}}
  \R{編者}{へん|しゃ} & 「\bkSeries」\R{編集委員会}{へん|しゅう|い|いん|かい} \\
   & \R{委員長}{い|いん|ちょう}　\R{小学校長}{しょう|がっ|こう|ちょう}（まとめ\R{役}{やく}） \\
   & \R{委員}{い|いん}　\R{数学}{すう|がく}・\R{電気工学}{でん|き|こう|がく}・\R{機械工学}{き|かい|こう|がく}・\R{教育学}{きょう|いく|がく}・\R{量子力学}{りょう|し|りき|がく}・\R{建築工学}{けん|ちく|こう|がく}・\\
   & \phantom{委員　}\R{半導体工学}{はん|どう|たい|こう|がく}の\R{大学}{だい|がく}\R{教授}{きょう|じゅ}、\R{幼稚園}{よう|ち|えん}\R{教諭}{きょう|ゆ}、\R{塾}{じゅく}\R{講師}{こう|し}、\\
   & \phantom{委員　}\R{特別支援学校}{とく|べつ|し|えん|がっ|こう}\R{教諭}{きょう|ゆ}、\R{心理士}{しん|り|し}、\R{科学}{か|がく}\R{解説者}{かい|せつ|しゃ}、\R{小学校}{しょう|がっ|こう}\R{教諭}{きょう|ゆ} \\
  デザイン & ユニバーサルデザイン\R{担当}{たん|とう}\R{委員}{い|いん} \\
  \R{挿絵}{さし|え} & イラスト\R{担当}{たん|とう}\R{委員}{い|いん} \\
  \R{版}{はん} & #3 \\
  \end{tabular}\par\vspace{3mm}
  {\color{bkVol0}\hrule height 1pt}\vspace{2mm}
  {\footnotesize #4}
  \end{minipage}
  \clearpage}
\makeatother
